(1) 均匀场近似中 $E(t)=V(t)/d$，电流密度 $\mathbf J=\sigma\mathbf E+\partial_t(\varepsilon\mathbf E)$。总电流 $I(t)=(A/d)[\sigma V_0\cos\omega t-\omega\varepsilon V_0\sin\omega t]$；前者同相，后者相差 $90^\circ$，不能在计算功率时直接将两电流幅度相加。
(2) 采用时间因子 $e^{i\omega t}$，复电流幅值为 $\tilde I=Y\tilde V$，其中 $Y=(A/d)(\sigma+i\omega\varepsilon)$，$Z=(d/A)/(\sigma+i\omega\varepsilon)$。导电与位移部分幅值比为 $\tan\delta=\sigma/(\omega\varepsilon)$，在介质工程中称损耗角正切，其正性来自 $\sigma>0$。
(3) 周期平均功率 $\langle P\rangle=\langle VI\rangle=(A/d)[\sigma V_0^2\langle\cos^2\omega t\rangle-\omega\varepsilon V_0^2\langle\sin\omega t\cos\omega t\rangle]=\sigma AV_0^2/(2d)$。一个周期 $2\pi/\omega$ 中耗散 $Q_{\rm cycle}=\pi\sigma AV_0^2/(\omega d)$，位移支路相应储能周期性变化，净贡献为零。
(4) 两种电流幅值相等时 $\sigma=\omega\varepsilon$，故 $\omega_c=\sigma/\varepsilon$；它是电荷弛豫率的倒数时间尺度。$\sigma\to0$ 时电流完全正交于电压，有功功率趋零；$\omega\to0$ 时位移电流趋零，介质表现为静态漏电电阻 $d/(\sigma A)$，不能将损耗正切发散解释为功率发散。
